\documentclass[12pt, letterpaper]{article}
\usepackage[margin=1in]{geometry}
\usepackage{amsmath}
\usepackage{amssymb}
\usepackage{mathtools}
\usepackage{comment}
\usepackage{enumerate}
\usepackage{changepage}
\usepackage{amsthm}
\usepackage{bbm}
\usepackage{tikz-cd}
\usepackage{xcolor}
\usepackage{hyperref}

\setlength{\parindent}{0pt}

%% code from mathabx.sty and mathabx.dcl
\DeclareFontFamily{U}{mathx}{\hyphenchar\font45}
\DeclareFontShape{U}{mathx}{m}{n}{
      <5> <6> <7> <8> <9> <10>
      <10.95> <12> <14.4> <17.28> <20.74> <24.88>
      mathx10
      }{}
\DeclareSymbolFont{mathx}{U}{mathx}{m}{n}
\DeclareFontSubstitution{U}{mathx}{m}{n}
\DeclareMathAccent{\widecheck}{0}{mathx}{"71}
\DeclareMathAccent{\wideparen}{0}{mathx}{"75}

\def\cs#1{\texttt{\char`\\#1}}

\allowdisplaybreaks

\DeclareMathOperator{\Id}{Id}
\DeclareMathOperator{\im}{im}
\renewcommand{\Re}{\text{Re}}
\renewcommand{\Im}{\text{Im}}
\newcommand{\D}{\mathbf{D}}
\newcommand{\0}{\mathbf{0}}
\newcommand{\1}{\mathbf{1}}
\newcommand{\loc}{\text{loc}}
\DeclareMathOperator{\dist}{dist}
\DeclareMathOperator{\Span}{Span}
\DeclareMathOperator{\sign}{sign}
\DeclareMathOperator{\supp}{supp}
\DeclareMathOperator{\op}{op}
\DeclareMathOperator{\tr}{tr}
\DeclareMathOperator{\x}{\mathbf{x}}
\DeclareMathOperator{\y}{\mathbf{y}}
\DeclareMathOperator{\z}{\mathbf{z}}

\title{Math 104 Homework 5}
\author{Due Date: Friday, July 31 at 11:59pm}
\date{}

\begin{document}

\maketitle

\textbf{Instructions:} Submit pdf solutions in LaTeX (preferred), or \textit{neatly} handwritten. Unless otherwise stated, you may cite any result from class or from Rudin's book, but all other assertions you make must be proved rigorously. At the top of your submission, please write down the names of the other students you collaborated with and any outside sources you consulted.

\vspace{0.5cm}

\begin{enumerate}
    \item 
    Let $(X, d_X)$ and $(Y, d_Y)$ be metric spaces. A function $f : (X, d_X) \to (Y, d_Y)$ is called \textit{Lipschitz continuous} if there is a constant $L > 0$ so that
    \begin{equation*}
        d_Y(f(p), f(q)) \leq Ld_X(p, q)
    \end{equation*}
    for all $p, q \in X$. For $\alpha > 0$, we say $f$ is \textit{H\"older continuous (with exponent $\alpha$)} if there is a constant $C > 0$ so that
    \begin{equation*}
        d_Y(f(p), f(q)) \leq C d_X(p, q)^\alpha
    \end{equation*}
    for all $p, q \in X$. In particular, H\"older continuous with exponent $1$ means the same thing as Lipschitz.
    \begin{enumerate}
        \item 
        Fix a point $p \in X$. Show that the map $g : X \to \mathbb{R}$ defined by $g(x) = d(x, p)$ is Lipschitz continuous, with constant $L = 1$.
        
        \item 
        Prove that any H\"older continuous function (with any exponent $\alpha > 0$) is uniformly continuous.

        \item 
        Prove that if $f$ is Lipschitz continuous and bounded, then $f$ is H\"older continuous with exponent $\alpha$, for any $0 < \alpha \leq 1$.

        \item 
        Let $[a, b]$ be a closed interval, and suppose $f : [a, b] \to \mathbb{R}$ is $C^1$ (meaning its first derivative exists and is continuous on $[a, b]$). Show that $f$ is Lipschitz continuous.

        \item 
        Suppose $f : \mathbb{R} \to \mathbb{R}$ is H\"older continuous with exponent $\alpha > 1$. Show that $f$ is constant.
        
    \end{enumerate}

    \vspace{0.5cm}

    \item 
    On the interval $(100, \infty)$, define functions $f(x) = x + \sin(x)$, $g(x) = x$, $p(x) = x + \sin(x)\cos(x)$, and $q(x) = p(x)e^{\sin(x)}$.
    \begin{enumerate}
        \item 
        Using facts we proved about limits and derivatives, prove that
        \begin{equation*}
            \lim_{x \to \infty} \frac{f(x)}{g(x)} = 1, \quad \lim_{x \to \infty} \frac{f'(x)}{g'(x)} \;\; \text{DNE}, \quad \lim_{x \to \infty} \frac{p(x)}{q(x)} \;\; \text{DNE}, \text{ and } \lim_{x \to \infty} \frac{p'(x)}{q'(x)} = 0.\\
        \end{equation*}
        \item 
        For both scenarios $\frac{f(x)}{g(x)}$ and $\frac{p(x)}{q(x)}$, explain why these results do not contradict L'H\^opital's rule.
        
    \end{enumerate}

    \textit{Note: even though we haven't learned about $\sin$ and $\cos$ yet, you are free to use any familiar properties of them you wish.}

    \vspace{0.5cm}

    \item 
    Let $f : (a, b) \to \mathbb{R}$ be continuous. Let $c \in (a, b)$, and suppose $f$ is differentiable on $(a, c)$ and $(c, b)$. Suppose that 
    \begin{equation*}
        \lim_{x \to c} f'(x)
    \end{equation*}
    exists and is a real number. Show that $f$ is differentiable at $c$. \textit{Hint:} Mean value theorem.
    

    \vspace{0.5cm}

    \item 
    Let $f : \mathbb{R} \to \mathbb{R}$ be $C^2$ (meaning its second derivative exists and is continuous on $\mathbb{R}$).
    \begin{enumerate}
        \item 
        Let $x \in \mathbb{R}$. For each $h > 0$, show there are numbers $\xi$ and $\eta$ with $x - h < \xi < x < \eta < x + h$ and
        \begin{equation*}
            \frac{f''(\xi) + f''(\eta)}{2} = \frac{f(x + h) + f(x - h) - 2f(x)}{h^2}.
        \end{equation*}
        
        \item 
        Show that, for each $x \in \mathbb{R}$,
        \begin{equation*}
            f''(x) = \lim_{h \to 0} \frac{f(x + h) + f(x - h) - 2f(x)}{h^2}.
        \end{equation*}

        \item 
        A function $f : \mathbb{R} \to \mathbb{R}$ is called \textit{convex} if for all $x, y \in \mathbb{R}$ and $0 \leq \lambda \leq 1$, we have
        \begin{equation*}
            f(\lambda x + (1 - \lambda)y) \leq \lambda f(x) + (1 - \lambda) f(y).
        \end{equation*}
        Geometrically, this means the graph of $f$ lies \textit{below} all of its secant lines. Show that any $C^2$ convex function $f$ satisfies $f''(x) \geq 0$ for all $x$. \textit{Hint:} convexity tells you something about the numerator in part (b).

        \item 
        Conversely, if a twice-differentiable function $f$ satisfies $f''(x) \geq 0$ for all $x$, show that $f$ is convex. \textit{Hint:} For $\lambda \in (0, 1)$ and $x < y \in \mathbb{R}$, the number $z \coloneqq \lambda x + (1 - \lambda)y$ satisfies $x < z < y$. Apply the mean value theorem on $[x, z]$ and $[z, y]$ to get an expression for $\lambda f(x) + (1 - \lambda) f(y)$.
        
    \end{enumerate}

    \vspace{0.5cm}

    \item 
    Let $f : [0, 1] \to \mathbb{R}$ be defined by
    \begin{equation*}
        f(x) = 
        \begin{cases}
            1, & \text{ if } x \in \mathbb{Q}\\
            0, & \text{ if } x \notin \mathbb{Q}.
        \end{cases}
    \end{equation*}
    Compute the upper and lower Riemann integrals of $f$. Is $f$ Riemann integrable?
    
    
    
    
\end{enumerate}



\end{document}     